% GENERATED FILE. Edit canonical lessons, then run npm run generate:books.
% canonical-source-hash: fnv1a64:c73501008594188b
% canonical-lessons: JA-W144-with-a-clock

\chapter{Writing: With a Clock}
\label{ch:ja-writing-with-a-clock}
\hlchaptermodality{144}

\begin{chapteropening}
\textbf{By the end of this chapter:} \emph{I can fill in a five-field form inside a three-minute clock, I know the 30-40 sign message this book's A1 paper adds, and I know that the JLPT has no writing paper to rehearse.}

Complete JA-W144-with-a-clock: five form fields against a three-minute timer, then at most one minute checking the small signs.
\end{chapteropening}

\section[Practice]{The same writing, with a clock}

\label{lesson:JA-W144-with-a-clock}



\begin{quote}
You wrote your own line with as long as you liked.

A task you can only do untimed is not yet a task you can do.
\end{quote}

\subsection*{Writing — timed production}
\textbf{What the real exam publishes, and what it does not.} The JLPT, the receptive test this track anchors to, has \textbf{no writing paper} at any level. It tests vocabulary and grammar, reading and listening, and nothing a candidate writes is scored. So there is no official A1 writing clock to rehearse against, and this lesson does not pretend there is one.

The task comes instead from the JF Standard's A1 Can-dos, which follow the CEFR: complete a small form, and write a short personal message. This track's assessment contract sets that message at \textbf{30 to 40 signs}, for a reader you name, such as a \textbf{\ja{ともだち}}. Japanese has no spaces to count words by, so length is counted in signs: every kana counts as one, small \textbf{\ja{っ}} and \textbf{\ja{ゅ}} included. Under thirty and the message is not done; well past forty and you spent minutes the task did not ask for.

Today you rehearse the form only, against \textbf{this book's own clock}, not a published one. Set a \textbf{three-minute timer}, then fill five fields: \textbf{\ja{なまえ}}, \textbf{\ja{とし}}, \textbf{\ja{しごと}}, \textbf{\ja{しゅみ}}, \textbf{\ja{かぞく}}. Fill every one. If a field needs a word this book has not given you, write a number or the nearest word you own; never look one up.

Speed is where stroke order shows. A sign written in the order you learned keeps its shape when you hurry; one written out of order is the one that falls apart. When the three-minute timer rings, \textbf{stop}, even mid-sign.

\begin{tcolorbox}[breakable,colback=teal!4,colframe=teal!35!black,title={Before you move on}]
Spend at most one minute checking, and do not rewrite the timed answer. Did you write \textbf{\ja{しゅみ}} with a small \textbf{\ja{ゅ}}? Under a clock, that is the sign a hurried hand writes full size.

Now the rule worth carrying out of this lesson. JLPT N5 passes on \textbf{80 of 180} overall, and a candidate also needs \textbf{38 of 120} on vocabulary, grammar and reading together and \textbf{19 of 60} on listening. Below either, the candidate fails, whatever the total says.

Writing appears nowhere in those numbers. A JLPT pass says nothing about what you just did, which is why this book scores writing as \textbf{its own paper}: 60 of 100, never averaged against the others.
\end{tcolorbox}
